%&LaTeX   (tabsize:4)
%%%%%%%%%%%%%%%%%%%%%%%%%%%%%%%%%%%%%%%%%%%%%%%%%%%%%%%%%%%%%%%%%%%%%%%%%%%%
% File:      latex_guide.tex
% Author:		Tom Scavo	<trscavo@syr.edu>
% Author:		and Vince Darley <vince@santafe.edu>
% Requires:  graphics.sty, alltt.sty, latex_defs.tex
%            html.sty, htmllist.sty
% Summary:   User's Guide to accompany the Alpha LaTeX macros 
%            in latex.tcl, version 4.0
% Note:      Be sure to typeset this file THREE TIMES!
%%%%%%%%%%%%%%%%%%%%%%%%%%%%%%%%%%%%%%%%%%%%%%%%%%%%%%%%%%%%%%%%%%%%%%%%%%%

\documentclass{report}
\usepackage{hyperref}
\usepackage{alltt}
\usepackage{graphics}

%%%%%%%%%%%%%%%%%%%%%%%%%%%%%%%%%%%%%%%%%%%%%%%%%%%%%%%%%%%%%%%%%%%%%%%%%%%
% File:      latex_defs.tex
% Author:    Tom Scavo  <trscavo@syr.edu>
% Requires:  html.sty, htmllist.sty
%%%%%%%%%%%%%%%%%%%%%%%%%%%%%%%%%%%%%%%%%%%%%%%%%%%%%%%%%%%%%%%%%%%%%%%%%%%

% Packages accompanying LaTeX2HTML:
%\usepackage{html,htmllist}
%\usepackage{htmllist}

% Why isn't this in html.sty?
% \newcommand{\latexhtml}[2]{\latex{#1}\html{#2}}

% URLs and e-mail addresses:
% \usepackage{url}  % defines a "smart" \url command
%\newcommand{\url}[1]{\texttt{#1}}

% Various logos:
\newcommand{\AmS}{{\protect\usefont{OMS}{cmsy}{m}{n}%
  A\kern-.1667em\lower.5ex\hbox{M}\kern-.125emS}}
% \newcommand{\BibTeX}{{\rm B\kern-.05em{\sc i\kern-.025em b}\kern-.08em
%     T\kern-.1667em\lower.7ex\hbox{E}\kern-.125emX}}
\newcommand{\BibTeX}{\textrm{B\kern-.05em\textsc{i\kern-.025em b}\kern-.08em
    T\kern-.1667em\lower.7ex\hbox{E}\kern-.125emX}}
\newcommand{\MakeIndex}{\textsl{MakeIndex}}
\newcommand{\Textures}{\textsl{Textures}}
\newcommand{\OzTeX}{O\kern-.03em z\kern-.15em\TeX}
\newcommand{\DirectTeX}{D\kern-.04em irect\kern-.10em\TeX}
\newcommand{\CMacTeX}{C\kern-.08em M\kern-.07em ac\kern-.10em\TeX}
% \newcommand{\TeXXeT}{T\kern-.1667em\lower.5ex\hbox{E}\kern-.125emX-\kern0pt%
%     -X\kern-.1667em\lower.5ex\hbox{E}\kern-.125emT}
% \newcommand{\emTeX}{e\kern-.03em m\kern-.15em\TeX}
% \font\logo=logo10 scaled\magstephalf 
% \newcommand{\MF}{{\logo META}\-{\logo FONT}}
\newcommand{\twoe}{{2${}_{\textstyle\varepsilon}$}}
\newcommand{\Alpha}{\textsl{Alpha}}
\newcommand{\AlphaLite}{\textsl{AlphaLite}}

% LaTeX environments:
\newcommand{\env}[1]{\texttt{#1}}

% LaTeX commands:
\newcommand{\cmd}[1]
	{\mbox{\texttt{\symbol{'134}#1}}}
\newcommand{\bracedarg}[1]
	{\symbol{'173}#1\symbol{'175}}  % {...}
\newcommand{\cmdarg}[2]{\texttt{\cmd{#1}\bracedarg{#2}}}
% \newcommand{\bracketedarg}[1]{\symbol{'133}#1\symbol{'135}}  % [...]
% \newcommand{\cmdoptarg}[3]
%   {\texttt{\cmd{#1}\bracketedarg{#2}\bracedarg{#3}}}
\newcommand{\ttstar}{\texttt{\symbol{'52}}}
% \newcommand{\starred}{\ttstar-ed}  % *-ed
% \newcommand{\anewline}{\cmd{\symbol{'134}}}  % \\
% \newcommand{\anewlinestar}{\cmd{\symbol{'134}\symbol{'52}}}  % \\*

% Alpha commands:
\newcommand{\acmd}[1]{\textsf{#1}}
% \newcommand{\entrylabel}[1]{\acmd{#1}\hfill}
% \newenvironment{commands}[1]
%   {\begin{list}{}
%   	{%
%   	\renewcommand{\makelabel}{\entrylabel}%
%     \settowidth{\labelwidth}{\acmd{#1}}%
%     \setlength{\labelsep}{1em}%
%     \setlength{\itemsep}{0em}%
%     \setlength{\parsep}{0em}%
%     \setlength{\rightmargin}{5em}%
%     \setlength{\leftmargin}{\labelwidth}%
%     \addtolength{\leftmargin}{\labelsep}%
%     \addtolength{\leftmargin}{2em}%
%     }%
%   }%
% 	{\end{list}}
\newenvironment{commands}
  {\begin{flushleft}\begin{tabular}{@{\hspace{1.5em}}ll@{}}}
	{\end{tabular}\end{flushleft}}
	
% Alpha menus:
\newcommand{\amenu}[1]{\textbf{\textsf{#1}}}

% Alpha keys:
\newcommand{\straightquote}{\texttt{\symbol{'015}}}  % '
\newcommand{\key}[1]{$\langle\hbox{\textsc{#1}}\rangle$}
%\begin{htmlonly}
%	\newcommand{\key}[1]{<\textsc{#1}>}
%\end{htmlonly}
% \newcommand{\key}[1]
% 	{\latex{$\langle$\mbox{\textsc{#1}}$\rangle$}\html{<\textsc{#1}>}}
% \newcommand{\key}[1]
% 	{\latex{$\langle\mbox{\textsc{#1}}\rangle$}\html{<\textsc{#1}>}}

% Files:
\newcommand{\filespec}[1]{\mbox{\texttt{#1}}}
\newcommand{\pkg}[1]{\filespec{#1}}
\newcommand{\filetype}[1]{\filespec{.#1}}

% Small verbatim stuff:
\makeatletter
\renewcommand{\verbatim@font}{\normalfont\ttfamily\small}
\makeatother


% Table of contents:
\setcounter{tocdepth}{3}

% Headers and footers:
\pagestyle{headings}

%%%%%%%%%%%%%%%%%%%%%%%%%%%%%%%%%%%%%%%%%%%%%%%%%%%%%%%%%%%%%%%%%%%%%%%%%%%
\begin{document}

\title{
	\includegraphics{logo.eps}  \\[6ex]
	User's Guide  \\
	\Alpha\ \LaTeX\ Macros \\
	{\large Version~5.0~early draft}
}
\author{
	Original manual by Tom Scavo \\
        Updates by the AlphaTcl community
}
\date{October 2003}
\maketitle

\tableofcontents

\chapter{Overview}

Welcome to the \Alpha\ \LaTeX\ macros, a library of Tcl code which
forms a mode for the text editor \Alpha, designed to ease the input and
processing of \LaTeX\ documents when using any of the Alpha family of
editors.

This document is still in need of a major update to document changes 
which have been implemented with versions 5.0 and above.

\section{Features}

\begin{itemize}
	\item automatically invokes \TeX\ mode when opening or saving a 
	\filetype{tex} document

	\item single-keystroke typesetting of \LaTeX\ documents (will even 
	typeset an untitled or unsaved window, or portions of windows)

    \item through the creation of a \TeX\ fileset, provides full
    support for `master documents' composed of many separate .tex files
    (which are `input' or `included' into the master document).  This
    support carries as far as automatically routing all
    document-centric commands (e.g. `typeset') from sub-documents to
    the master document.
    
	\item with Alpha 7.6, Alpha 8 (on MacOS classic), and Alpha X
	or Alphatk (on MacOS X), it works with all known Macintosh \TeX\
	implementations, including \OzTeX, \Textures, \CMacTeX,
	Euro-\OzTeX, \DirectTeX, \DirectTeX~Pro, TeXshop, iTeXMac.

    \item with Alphatk on any platform (Windows, Unix,
    MacOS X), it works with all known command-line TeX tools
    including teTeX, MikTeX, Yap, Xdvi, including the tools
    running under `cygwin' on Windows.

    \item with Alphatk or Alpha X it can act as an interactive shell
    for your \TeX\ processing.
    
    \item where the TeX system supports it, synchronising back
    and forth between the source .tex document and a dvi viewer
    is provided (\Textures, Yap, Xdvi, at least).

	\item additional support provided for \Textures\ version~1.8

	\item dynamic menus and menu items
	
	\item choice of long or short \LaTeX\ menu; optional floating menu, 
	as well

	\item color syntax highlighting of \LaTeX\ keywords

	\item intelligent treatment of highlighted text (called ``wrapping'')

	\item easily creates \LaTeX\ document templates, complete with 
	indentation and tab stops

	\item handy text-to-\LaTeX\ conversion utilities

	\item quickly navigate and select commands, environments, 
	subsections, sections, and chapters

	\item ``smart quotes'' and ``smart dots'' with on-the-fly escape; 
	``smart'' subscripts and superscripts as well

	\item command-double-click feature that chases references and 
	citations, or opens \cmd{input} and \cmd{include} files; also 
	opens \cmd{bibliography}, \cmd{includegraphics}, \cmd{usepackage}, 
	and \cmd{documentclass} files

	\item a handy mark menu for navigating large \LaTeX\ documents

	\item follows closely the terminology and organization of Leslie 
	Lamport's \textsl{\LaTeX: A Document Preparation System} [Reading, 
	MA: Addison-Wesley, 1985, 1994 (ISBN 0-201-52983-1)]
	
    \item closely integrated with new features of AlphaTcl 7.x and
    8.x (for example the user can easily add menu items to most TeX
    menus, and can set user-defined key-bindings to most menu
    items)
	
	\item sophisticated support for Electric Completions and Expansions
	packages.
	
	\item smart indentation of pasted text if the smartPaste package is 
	activated.
\end{itemize}

\section{Documentation}
\label{sec:docs}

Pull down the System help menu and choose the command 
\acmd{LaTeX Help}.  This will open a new window with a brief 
introduction to the \Alpha\ \LaTeX\ macros.  From this window, you can 
access numerous other documentation files using \Alpha's built-in 
hypertext capability.  Just click on any of the following links in 
\acmd{LaTeX Help}:
\begin{flushleft}
	\begin{tabular}{@{}l@{}}
		\textbf{User's Guide:}  an introduction to the \Alpha\ \LaTeX\ macros 
		(this file)  \\
		\textbf{\LaTeX\ Menus:}  commands and bindings (organized by menu)  \\
		\textbf{\LaTeX\ Key Bindings:} commands and bindings (organized by 
		command key)  \\
	\end{tabular}
\end{flushleft}
Note that the User's Guide is available in both \LaTeX\ and \textsc{html} 
formats.  (The latter is online at URL
\begin{flushleft}
    \begin{tabular}{@{}l@{}}
        \url{no-current-url}
    \end{tabular}
\end{flushleft}
and was created by the program \filespec{TeX4ht}.

In addition to the above \Alpha\ \LaTeX\ documents, there are also the 
more general Web pages
\begin{flushleft}
    \begin{tabular}{@{}l@{}}
        \textbf{An Introduction to \LaTeX}  \\
        \url{http://www.latex-project.org/intro.html} \\
        \textbf{An Introduction to \LaTeX\ and \AmS-\LaTeX}  \\
        \url{http://where is this}
    \end{tabular}
\end{flushleft}

\section{\LaTeXe}

The \Alpha\ \LaTeX\ macros support \LaTeXe.  \Alpha\ no longer
directly supports
the 2.09 macros (although many are unchanged in \LaTeXe so if you must
you could probably still manage).  \LaTeXe\ will typeset a 2.09
document automatically, using what is called ``compatibility mode''.
Most of the \LaTeX~2.09 commands and environments have been preserved
in \twoe, making the transition from 2.09 to \twoe\ relatively painless
(from the user's point of view, at least).

\section{Installation}

\Alpha\ is configured to use the \LaTeX\ macros right out of the box, 
so there is no installation process per se.  However, there are a 
number of flags and variables that control the inner workings of 
\filespec{latex.tcl} that may be changed at the user's discretion.  

The following \TeX\ mode flags and variables may be accessed by 
pulling down the \amenu{Config} menu and opening the 
\acmd{Preferences} dialog on the \amenu{Current Mode} submenu.  See 
the \acmd{Alpha Manual} (on the System help menu under the question 
mark) for more information about \Alpha's global preferences.

\subsection{Flags}
\label{sec:flags}

% {\renewcommand*{\descriptionlabel}[1]{\hspace\labelsep \normalfont #1}
\begin{description}
	\item[\acmd{buildPkgsSubmenu}] The \amenu{Packages} submenu is an 
	optional submenu containing a list of all \filetype{tex} and 
	\filetype{sty} files in your \TeX\ search path.  Choosing a filename 
	from the list inserts the corresponding \cmd{usepackage} command 
	into the preamble of the current document.  By default, however, the 
	\amenu{Packages} submenu is not built when the \LaTeX\ macro package 
	is loaded.  To build this submenu on-the-fly, enable the flag 
	\acmd{buildPkgsSubmenu} as described above, and then choose 
	\acmd{Rebuild Documents Submenu} on the \amenu{Documents} submenu (see 
	section~\ref{sec:documentsSubmenu}).  Thereafter, the 
	\amenu{Packages} submenu will be built automatically along with the 
	\LaTeX\ menu.

	\item[\acmd{deleteObjNoisily}] One of the basic \filespec{latex.tcl} 
	operations is to insert an object into the current document.  If, at 
	the time the insertion command is issued, there is a selection 
	(i.e., text is highlighted), then the program behaves differently 
	depending on the value of the flag \acmd{deleteObjNoisily}.  If set 
	to true, the user will be prompted before any selected text is 
	deleted.  If, on the other hand, this flag is false, then the 
	selection is replaced quietly and without warning (although it may 
	be undone).  By default, \acmd{deleteObjNoisily} is set to true.  
	NOTE: Not all objects are ``inserted'' into the document since 
	sometimes there is an attempt to ``wrap'' the current selection.  
	See section~\ref{sec:basicOps} for more information.

	\item[\acmd{deleteEnvNoisily}] Before an environment is inserted 
	into the document, the program checks to see if there is a 
	selection.  If so, and the flag \acmd{deleteEnvNoisily} is set to 
	true, the user is asked whether or not the current selection should 
	be replaced; if false, the current selection is deleted without 
	warning.  Note that the default value of \acmd{deleteEnvNoisily} has 
	been set to true.  Like objects, environments may wrap, so sometimes 
	the current selection is treated differently.  See 
	section~\ref{sec:basicOps} for details.

	\item[\acmd{promptNoisily}] Some environment commands prompt the user for
	input.  As mentioned below, if \acmd{useStatusBar} is set to true, the
	prompt is displayed on the thin status bar at the top (or bottom) of the
	screen.  This is less obtrusive than a dialog, but may go unnoticed at
	first, so if \acmd{promptNoisily} is set to true (which it is, by
	default) and \acmd{useStatusBar} is enabled, the program beeps prior to
	displaying the prompt.  You can turn off this annoying sound by invoking
	the \acmd{Flags} command on the \amenu{Current Mode} submenu of the
	\amenu{Config} menu and removing the check on \acmd{promptNoisily}.

	\item[\acmd{runTeXInBack}] If true, typesetting will occur in the 
	background.  This flag is false by default.  The main 
	\amenu{Process} menu contains options for both foreground and 
	background typesetting.  This preference setting simply defines 
	which of the two is bound to a \key{Cmd-T} keypress, and which to 
	\key{Shift-Option-T}.

	\item[\acmd{searchNoisily}] Many commands cause 
	\filespec{latex.tcl} to search the current document.  If a 
	search fails, and \acmd{searchNoisily} is set to true, the 
	program displays a message on the status bar and beeps.  If, on 
	the other hand, \acmd{searchNoisily} is set to false, only 
	the message is displayed.  By default, \acmd{searchNoisily} 
	is set to true.

	\item[\acmd{smartDots}] By default, \filespec{latex.tcl} 
	replaces three consecutively typed dots (\texttt{...}) with the 
	\LaTeX\ command \cmd{ldots}.  To escape the effect of 
	\acmd{smartDots}, press the \key{Delete} key on-the-fly.

	\item[\acmd{smartQuotes}] If this flag is set to true, pressing the
	single quote key \key{\straightquote} will generate \verb/`/ or \verb/'/
	automatically depending on the context.  Similarly, pressing the double
	quote key \key{\"} generates \verb/``/or \verb/''/, whichever is
	required.  Set \acmd{smartQuotes} to false if you want the single and
	double quote keys to insert \key{\straightquote} and \key{\"} literally,
	or press the \key{Delete} key to escape the effect of \acmd{smartQuotes}
	on-the-fly.

	\item[\acmd{smartScripts}] When this flag is enabled (which it is by 
	default), the \verb|^| and \verb|_| keys on a U.S. keyboard are 
	bound to the commands \acmd{superscript} and \acmd{subscript}, 
	respectively, on the \amenu{Formulas} submenu of the \LaTeX\ menu 
	(see section~\ref{sec:formulasSubmenu}).  Press the \key{Delete} key 
	to escape the effect of \acmd{smartScripts} on-the-fly.

% 	\item[\acmd{useBoxMacro}] If true, the body of a figure environment 
% 	includes a call to the box-making macro stored in the variable 
% 	\acmd{boxMacroName} (see section~\ref{sec:variables} below); if 
% 	false, the body of a figure environment is left blank, unless there 
% 	is a current selection, in which case the selection is wrapped.

	\item[\acmd{useBrackets}] In \LaTeX, the \env{displaymath} 
	environment is equivalent to \cmd{[}\ldots\cmd{]}.  If you prefer to 
	use the latter, set \acmd{useBrackets} to true.  By default, 
	\acmd{useBrackets} is set to false, that is, the \env{displaymath} 
	environment is used to construct multi-line math displays.  Note: By 
	default, \filespec{latex.tcl} \emph{always} uses 
	\cmd{[}\ldots\cmd{]} inline (unless \acmd{useDollarSigns} is set to 
	true---see below).  This setting applies whenever \LaTeX mode 
	inserts math environments (usually due to a menu selection).

	\item[\acmd{useDollarSigns}] Support is provided for both the \TeX\ 
	and \LaTeX\ methods of invoking inline math mode (see the 
	\amenu{Math Modes} submenu in section~ \ref{sec:mathModeSubmenu} for 
	the various options), but only one of these is bound to command keys 
	(namely, \key{Ctl Cmd M} and \key{Ctl Opt Cmd M}, by default).  This 
	is what the flag \acmd{useDollarSigns} does.  If set to true, 
	\filespec{latex.tcl} uses dollar signs to delimit inline math mode 
	(\verb|$|\ldots\verb|$| and \verb|$$|\ldots\verb|$$|), whereas if it 
	is false, \LaTeX\ notation will be used (\cmd{(}\ldots\cmd{)} and 
	\cmd{[}\ldots\cmd{]}).  By default, \acmd{useDollarSigns} is set to 
	false---the \LaTeX\ way of doing things.

	\item[\acmd{useStatusBar}] This flag determines whether or 
	not the status bar is used when prompting for user input.  (The 
	status bar is a long, thin message area at the bottom of your 
	screen.)  Use of \Alpha's status bar is enabled in \TeX\ mode, by 
	default.  See the related flag \acmd{promptNoisily} above.
\end{description}%}
 
\subsection{Variables}
\label{sec:variables}

% {\renewcommand*{\descriptionlabel}[1]{\hspace\labelsep \normalfont #1}
\begin{description}
	\item[\acmd{boxMacroNames}] This \TeX\ mode variable contains a list 
	of names of box-making macros used in the body of a figure 
	environment.  The standard \LaTeX\ commands \cmd{includegraphics} 
	and \cmd{includegraphics\ttstar} (both part of the \LaTeXe\ 
	\pkg{graphics} package) are included in this list by default.

	\item[\acmd{citeCommands}] Any command listed as a ``cite'' command 
	is command-double-clickable (see section~\ref{sec:dblClicking}).  
	The standard \LaTeX\ commands \cmd{cite} and \cmd{nocite} are 
	included in this list by default.

	\item[\acmd{fillcolumn}, \acmd{leftFillColumn}] See the \acmd{Alpha
	Manual} on the menubar's Help menu for more information about these
	variables.

	\item[\acmd{funcExpr}] In \TeX\ mode, \acmd{funcExpr} is a regular 
	expression used to search for a subsection header (see the 
	commands \acmd{Next Subsection} and \acmd{Prev Subsection} 
	described in section~\ref{sec:gotoSubmenu}) and to build the 
	\amenu{funcs} pop-up menu (see section~\ref{sec:funcs}).

	\item[\acmd{funcExprAlt}] In \TeX\ mode, \acmd{funcExprAlt} is a 
	regular expression used to search for a section header (see the 
	commands \acmd{Next Section} and \acmd{Prev Section} described in 
	section~\ref{sec:gotoSubmenu}).

	\item[\acmd{lineWrap}] If this variable is set to \acmd{Auto}, the
	program automatically inserts a carriage return as the cursor
	nears the end of a line (the length of which is defined by the
	variable \acmd{fillColumn} described in
	section~\ref{sec:variables} below); otherwise, the line extends
	indefinitely to the right (until the \key{Return} key is pressed,
	of course).  By default, \acmd{lineWrap} is turned on in \TeX\
	mode.  See the \acmd{Alpha Manual} on the menubar's Help menu for
	more information.

	\item[\acmd{parseExpr}] The variable \acmd{parseExpr} is a regular 
	expression used to build the \amenu{funcs} pop-up menu.  See 
	section~\ref{sec:funcs} for more information.

	\item[\acmd{prefixString}] This variable is used in 
	conjunction with the \acmd{Comment Line} command on \Alpha's 
	\amenu{Text} menu.  In \TeX\ mode, this string is set to 
	``\verb*|% |'' by default, which makes \acmd{comment line} very 
	useful for commenting out large blocks of \LaTeX\ code.

	\item[\acmd{refCommands}] Any command listed as a ``ref'' command is 
	command-double-clickable.  The standard \LaTeX\ commands \cmd{ref} 
	and \cmd{pageref} are included in this list by default.

	\item[\acmd{TeXSearchPath}] Whenever \LaTeX\ mode searches for 
	files, it examines each directory in this user-defined list.  Such a 
	search is normally triggered by the user command-clicking on a 
	citation, reference or input command in a \LaTeX\ document.  In such 
	cases, \Alpha\ 
	attempts to find the appropriate document and open it for the user.
	
	\item[\acmd{wordBreak}] This variable holds a regular expression
	that define a ``word'' in \TeX\ mode.  (A ``word'' is any text
	string that is double-clickable.)  See the \acmd{Alpha Manual} in
	the help menu for more information.

%	\item[\acmd{wrapBreak}, \acmd{wrapBreakPreface}] These variables are 
%	similar to \acmd{wordBreak} and \acmd{wordBreakPreface} above, 
%	except that they are used by \Alpha\ to wrap lines, not delineate 
%	words.  See the \acmd{Alpha Manual} on the System help menu (under 
%	the question mark) for more information.
\end{description}%}

\subsection{Tips and tricks}

A useful installation trick that you might want to put in your 
\filespec{prefs.tcl} file (opened by choosing \acmd{Edit Prefs} 
on the \amenu{Global} submenu of the \amenu{Config} menu) 
is the following key binding:
\begin{alltt}
   Bind \straightquote{}n\straightquote <cs> \{TeX::newLaTeXDocument\}
\end{alltt}
With this binding, it's easy to bring up a new \LaTeX\ document no 
matter where you are or what you're doing.  Regardless of the file 
you're currently editing, simply press \key{Shf Cmd N} to open a new 
\LaTeX\ document (see the \acmd{New Document} command in 
section~\ref{sec:documentsSubmenu} for details).  If you use
\Alpha's `New Document' package, this capability is automatically
part of the standard `New Document' dialog.

\Alpha\ is completely customizable, but it's not a good idea to 
modify its Tcl files directly.  Instead, put your modifications in 
preferences files designed specifically for this purpose.  Besides 
the global preferences file \filespec{prefs.tcl} mentioned above, 
there are also mode-specific preferences files.  \TeX\ mode, for 
instance, has its own prefs file called \filespec{TeXPrefs.tcl}.  To 
edit this prefs file, open a \filetype{tex} file and choose the 
command \acmd{Edit Prefs File} on the \amenu{TeX Mode Prefs} submenu of the 
\amenu{Config} menu.  All \TeX-related modifications should be 
placed in this preferences file.

Note: All preferences files are stored separate to \Alpha\ and are not
touched when you upgrade your version of \Alpha.

The \LaTeX\ menu can be modified to suit your needs.  Three different
methods allow you to insert new commands into any of the \amenu{Text}
or \amenu{Math} submenus, presented here in increasing order of 
difficulty.

The first method will define a simple \LaTeX\ ``command'' item that is
followed by braces.  Select the \acmd{Add Menu Template Item} from the
\amenu{LaTeX Menu} submenu.  Choose a menu that should contain the new
item.  A new dialog appears with three text fields.  In the first field
enter the name of the item as you want it to be presented in the menu.
The second field contains text looking like 
\textbf{\\XXX\textbullet\{\textbullet} --- replace \textbf{XXX} with 
the name of the command as it should appear in the window when the 
menu item is called.  If you want to include any default text within 
the braces that will follow this item, you can insert that in either 
the second or the third text fields.  Click on the ``OK'' button to 
save the changes.  All menu items added in this way can be changed 
later using the \acmd{Edit Template Item} command, or deleted using the 
\acmd{Delete Template Item} command.

The second method requires a little working knowledge of AlphaTcl.
Suppose, for example, that you want to insert a new \acmd{My Template}
item to the \amenu{Arrows} menu.  Open your \filespec{TeXPrefs.tcl}
file by selecting the \acmd{Edit Prefs File} item in the \amenu{Config
\textgreater TeX Mode Prefs} submenu.  Paste in a command that looks
like this:

\begin{alltt}
	menu::insert "Arrows" items 0 "My Template"
\end{alltt}

Now you need a Tcl procedure that will be evaluated when this item is
selected.  It should be named \textbf{TeX::MyTemplate}, i.e. scrunch the
name together by removing all spaces, but retain upper and lower cases
to match the name of the menu item.  Place this new procedure in your
\filespec{TeXPrefs.tcl} file, and ``evaluate'' this file by pressing
\key{Command-L}.  The new item should now appear in the menu.  If it
does, save the preferences file, and close it.  This file will be
``sourced'' automatically the next time that Alpha is started and a
.tex file has been opened.

The third method requires more extensive knowledge of AlphaTcl.  You 
must turn on the package ``Smarter Source'', and place a copy of the 
file \filespec{latexMenus.tcl} in your ``Smarter Source'' folder.  Look 
for the definition of the menu in question, and modify it to include 
the items that you want.  Restart Alpha, and your customized \LaTeX\ 
menu should include the items you added.  Unfortunately, you're not 
done yet.  You must now modify the ``menu procedure'' that is called by 
the given menu in order for your new item to actually do anything!  
This procedure is listed in the definition of the menu, and if it is 
not in the file \filespec{latexMenus.tcl} then you'll have to 
Command-Double-Click on the name of the procedure to find it, and then 
copy \emph{that} file to your ``Smarter Source'' folder as well, and 
then modify it accordingly.

An alternative to this third method would be to simply use the
``Macros'' package to define a new macro, edit it to do what you want,
and then assign it a keyboard shortcut so that it can be called
whenever you want.  Not as elegant, perhaps, but in the end just as 
effective and less complicated.

Feel free to ask for more help in the AlphaTcl mailing lists.

\section{\TeX\ Filesets}

A very important capability of \Alpha's \LaTeX\ mode is the seamless 
handling of documents composed of many separate files on disk. 
Through the one-off action of creating a \TeX-fileset you inform
\Alpha\ of the relationship between these parts and it automatically
uses that information from that point onwards.  This means, for
example, when you attempt to typeset a sub-document, \Alpha\
notices, retrieves the master document instead and processes it for
typesetting.  Here is a partial list of the actions which
are automatically aware of multi-part documents:
\begin{enumerate}
    \item Typesetting, \BibTeX'ing and all other actions in the
    \amenu{Process} menu.

    \item Double-clicking on references to jump to their matching
    label.

    \item Searching (the standard search dialog provides option to
    search all files in any defined fileset).

    \item Synchronicity features (.tex to .dvi and .dvi to .tex
    synchronisation).

\end{enumerate}

In addition all parts of the multi-part TeX/LaTeX document are listed
in \Alpha's \amenu{Fileset} menu which therefore provides quick and easy
access to editing these files.

To create a \TeX\ fileset, select \amenu{Fileset \textgreater Utilities
\textgreater New Fileset}.  When prompted for the type of the new
fileset, make sure you select \TeX.

There are, however, some known issues involving the proper resolution 
of \cmd{include} file path names when they don't exist in the same 
directory as your base file.  You can specify \"Additional Related 
Files\" for your \TeX fileset to ensure that these are included in 
the list of files.  Whenever these files are the active \Alpha 
window, typesetting and other \amenu{Process} commands will still 
operate on the base \TeX file.

\section{Basic operations}
\label{sec:basicOps}

The \Alpha\ \LaTeX\ macros revolve around two basic operations called 
\acmd{insertObject} and \acmd{wrapObject}.  Basically, 
\acmd{insertObject} is a call to the primitive procedure 
\acmd{insertText} preceded by the automatic deletion of previously 
selected text (this behavior is easily changed, however, by resetting 
the flag \acmd{deleteObjNoisily} described in 
section~\ref{sec:flags}).  For example, if there is no current 
selection, choosing the command \acmd{alpha} from the \amenu{Greek} 
submenu of the \LaTeX\ menu inserts the \LaTeX\ command \cmd{alpha} at 
the insertion point; otherwise, if there is a selection, it is 
replaced with the string ``\cmd{alpha}''.  In other words, 
\acmd{insertObject} works just like the familiar \acmd{Paste} command 
on the \amenu{Edit} menu.  It turns out that a large number of 
commands in \filespec{latex.tcl} rely on \acmd{insertObject}, but 
sometimes it's faster to type the desired \LaTeX\ command directly.  
(Even faster is to use the corresponding command keys, but more on 
that later.)  On the other hand, if you forget the syntax of a 
particular \LaTeX\ command, it's sometimes easier to look it up on the 
\LaTeX\ menu than it is in a reference manual.

The complementary operation to \acmd{insertObject} is called 
\acmd{wrapObject}.  The difference between the two is the way the 
latter treats the current selection, that is, \acmd{wrapObject} 
inserts its argument at the insertion point (just like 
\acmd{insertObject}), but if there is a selection, \acmd{wrapObject} 
cuts it out (without effecting the state of the Clipboard) and inserts 
it in the middle of the chosen command.  For example, consider the 
\LaTeX\ menu command \acmd{footnote} on the \amenu{Miscellaneous} 
submenu (see section~\ref{sec:miscSubmenu}).  This command inserts the 
string ``\cmdarg{footnote}{}$\bullet$'' (without the double quotes, of 
course) into the document, positioning the insertion point between the 
pair of braces.  The user then types the text to be footnoted and 
presses the \key{Tab} key, after which the tab stop macro finds (and 
deletes) the bullet $\bullet$ at the end of the string.  (Note: Use 
\key{Opt Tab} to insert a literal tab character into the document.)  
On the other hand, if a selection exists at the time the 
\acmd{footnote} command is issued, the selection itself is surrounded 
by \LaTeX's \cmd{footnote} command, and the insertion point is brought 
to the end of the selection automatically.  Some commands, for better 
or worse, even go so far as to insert the selection into one of 
several competing positions within the command string.  The 
\acmd{fraction} command on the \amenu{Formulas} submenu 
(section~\ref{sec:formulasSubmenu}) is a good example of this type of 
behavior.  It assumes the current selection (if there is one) is the 
numerator of the fraction to be typeset, cutting and pasting 
accordingly.

The concept of wrapping is carried one step further in the case of 
environments.  Suppose you want to center an existing \env{tabular} 
environment, for example.  Just select the \env{tabular} environment 
to be centered and choose the \acmd{center} command from the 
\amenu{Environments} submenu on the \LaTeX\ menu (see 
section~\ref{sec:envSubmenu}).  The resulting \env{center} environment 
will completely surround the existing \env{tabular} environment, 
indenting the latter one tab stop to the right.

Not all environments wrap, however.  Those environments whose body is 
very structured (such as \env{enumerate}, \env{itemize}, 
\env{description}, \env{thebibliography}, \env{tabular}, \env{array}, 
\env{eqnarray}, and \env{eqnarray*}) do not.  Instead, these 
environments simply insert text into the document.  If there happens 
to be a selection at the time one of these commands is issued, an 
alert appears asking if the selection should be deleted.  To turn this 
alert off, simply toggle the flag \acmd{deleteEnvNoisily} (see 
section~\ref{sec:flags}) in the \acmd{Flags} dialog on the 
\amenu{Current Mode} submenu on the \amenu{Config} menu.


\section{Whitespace}

Before continuing, let me say a few words about whitespace.  In 
virtually all cases, superfluous whitespace in command strings has 
been deleted.  For example, objects inserted with \acmd{insertObject} 
(a sizable portion of \filespec{latex.tcl}'s functionality) do not 
routinely insert a trailing space character.  Instead, the user must 
decide whether or not space should immediately follow a particular 
\LaTeX\ control word, since sometimes it's needed and sometimes it's 
not.  Disabling the \acmd{indentOnReturn} preference for \LaTeX\ mode also 
prevents a lot of extraneous whitespace from being inserted into your 
document\ldots but then it won't look so good!

\section{Customising}

\LaTeX\ mode contains a vast array of preferences which can be rather 
confusing.  Here are some not so obvious ways to customise it.

Always opening the log file? Add a key-binding (here with 
\key{Opt Cmd O}):

\begin{verbatim}
Bind `o' <co> {TeX::mp::Processmenu "" "Open .log"} TeX
\end{verbatim}

Want to add some optional parameters to an automatically inserted
latex environment? (E.g. you like \env{enumerate} with a parameter in 
square brackets `[]'), then add the following line to your 
\filespec{TeXPrefs.tcl} (where the stars are bullets):

\begin{verbatim}
set TeXbodyOptions(enumerate) "\[*a|i*\]"
\end{verbatim}


\section{Bugs}
\label{sec:bugs}

Comments, suggestions, and bug reports are certainly welcome.  In 
fact, many of the improvements and features in this version of 
\filespec{latex.tcl} were suggested by \Alpha\ \LaTeX\ users.  (Some of 
them even sent me code!).

Alpha is privileged by an active and helpful user community whose main
communication channels are two mailing lists you are invited to subscribe 
to.  The developers' list is for discussing technical issues, while the users' list
is for all sorts of questions and hints --- it is read also by the developers.
Questions are usually answered within a day, and often they lead to
improvements in subsequent versions of Alpha and its documentation.
Go to:

\url{http://lists.sourceforge.net/lists/listinfo/alphatcl-users}
    
\url{http://lists.sourceforge.net/lists/listinfo/alphatcl-developers}

In addition there is the bugzilla bug tracking system which should be
used to report bugs: 

\url{http://www.purl.org/net/alpha/bugzilla/}


The following are known bugs in \filespec{latex.tcl}:
\begin{itemize}
	\item When wrapping, the \acmd{frac} command does not remove redundant 
	parentheses.

	\item The \acmd{options} commands assumes the \cmd{documentclass} 
	command already has an optional parameter (which it does if the 
	document template was inserted via the \LaTeX\ menu).  Moreover, the 
	\acmd{options} command does not check for duplicate options.

	\item The commands \acmd{Prev Command Select With Args} and 
	\acmd{Next Command Select With Args} will not select \LaTeX\ commands 
	whose argument(s) contain braces.

	\item  The \amenu{Goto} submenu could be better organized (and will 
	be, if I can ever think of a good set of command keys!).

	\item The results of \acmd{Delete Comments} can not be undone (but it 
	works!).
	
	\item  The \acmd{Any TeX File} command on the \amenu{Process} 
	submenu does not open the correct folder.
	
	\item  The at-symbol @ is not recognized as a \LaTeX\ command 
	character in \filetype{sty} files.
	
	\item The command keys for \acmd{subscript} and \acmd{superscript} 
	are not compatible with international keyboards.  However they can be
	set by the user manually.
	
	\item  The marking algorithm should ignore comments.
	
	\item The \amenu{Process} submenu gets confused if the current file 
	has a name containing meta-characters.  The only known fix is to 
	avoid the characters \verb|<|, \verb|(|, \verb|;|, \verb|!|, 
	\verb|^|, and \verb|/| in filenames.
\end{itemize}


\section{Acknowledgments}

Tom Scavo wrote most of \LaTeX\ mode and this documentation.  The 
present author, Vince Darley can take very little credit.

Here are Tom's original acknowledgements.

Numerous people have made significant contributions to the \Alpha\ 
\LaTeX\ macros.  You will find their names and initials scattered 
throughout this and other \Alpha\ documents.  
Tom~Pollard 
\href{mailto:pollard@cucbs.chem.columbia.edu}{pollard@cucbs.chem.columbia.edu}
and Vince~Darley 
\href{mailto:vince@santafe.edu}{vince@santafe.edu}
have been especially helpful and deserve a lot of credit.  Of course, 
none of this would have been possible without the support and 
encouragement of \Alpha's author,
\begin{flushleft}
    \begin{tabular}{@{\hspace{1.5em}}ll@{}}
        Pete~Keleher 
        & 
        \href{mailto:keleher@cs.umd.edu}{keleher@cs.umd.edu}
    \end{tabular}
\end{flushleft}
whom I heartily thank.

\chapter{Menus}

\section{The \LaTeX\ menu}

Upon entering \TeX\ mode, either 
\begin{enumerate}
	\item \emph{manually} (by choosing \acmd{TeX} from the pop-up mode 
	menu on the status bar at the bottom of your screen); or

	\item \emph{automatically} (whenever a \filetype{tex} or 
	\filetype{sty} file is opened or saved);
\end{enumerate}
a new menu appears in the menu bar.  The \LaTeX\ menu provides access 
to scores of procedures loaded automatically the first time \TeX\ mode 
is entered.

The bulk of the \LaTeX\ menu contains items for marking up ``Text'' and
``Math'' items.  You can condense the submenus in these two categories by
selecting the \acmd{Compress Text/Math Menus} command in the \amenu{LaTeX Menus}
submenu.  All of the commands will be available even if the menus have been 
compressed, these preferences only affect the path that your mouse will have 
to take in order to access them.
\label{pg:compressMenus}

The \LaTeX\ menu follows closely the organization and terminology of 
Lamport's \textsl{\LaTeX: A Document Preparation System} [second 
edition, Addison-Wesley, 1994], especially chapter~3.  Many people 
agree that the \textsl{\LaTeX} book is still the definitive \LaTeX\ 
reference.  In conjunction with \textsl{The \LaTeX\ Companion} by 
Goossens, Mittlebach, and Samarin [Addison-Wesley, 1994 (ISBN 
0-201-54199-8)], these two books constitute the ``official'' \LaTeXe\ 
documentation.  These books, as well as Knuth's classic \textsl{\TeX 
book} [Addison-Wesley, 1986 (ISBN 0-201-13448-9)], should be on every 
serious \LaTeX\ user's desk.

The \LaTeX\ menu is organized into four parts: general commands, 
document-related commands, paragraph mode commands (that is, text 
commands), and math mode commands.  Each group of commands is 
separated by a thin grey line on the \LaTeX\ menu.  The order of the 
commands on any given submenu is significant insofar as possible.  For 
example, the various commands on the \amenu{Environments} submenu 
mirror the corresponding command keys, while other submenus follow the 
ordering found in the \textsl{\LaTeX} book.  We'll try to point out 
these organizational aids as we go along.

A brief description of each available command follows.  See  
section~\ref{sec:docs} for pointers to other help documents.

\subsection{Menu Keyboard Shortcuts}

Most of the \LaTeX\ menu items can be accessed by keyboard shortcuts that are
specific to \TeX\ mode.  Many of these shortcuts can be changed to any
combination that might seem more intuitive.  

Before you change a menu item keyboard shortcut, you should probably make sure
that the new modifier/keystroke combination is not already used by some other
function.  Select the \acmd{Describe Binding} menu command in the
\amenu{Config} menu to determine if this is the case.

Once you have decided upon some keyboard shortcuts that you want to change,
use the \acmd{Assign Menu Bindings} command in the \amenu{LaTeX Menu} submenu
to first select the relevant menu, and then change the bindings that appear in
the dialog by clicking on the relevant ``Set'' buttons.  When you have
finished, click on ``OK'', and then either choose another menu or select
``Finish'' and press ``OK''.  All changes will take effect immediately, and
will be remembered between editing sessions.

\subsection{General commands}

The \Alpha\ \LaTeX\ mode can interact with a very large number of
other applications for typesetting, viewing, converting documents, etc.

In particular recent versions of Alphatk can be used as external
editors for command-line \LaTeX\ tools.

Most of this interactions occurs through the \amenu{Process} submenu.

\subsubsection{\amenu{Process} submenu}

\begin{commands}
	\acmd{Typeset file.tex} & \key{Cmd T} \\
	\acmd{View file.dvi} & \key{Shf Cmd V} \\
	\acmd{Print file.dvi} & \key{Shf Cmd P}
\end{commands}
If you use \Textures, \OzTeX, \CMacTeX, or \DirectTeX, or any
command-line tool on MacOS X, Windows or Unix, you'll be happy 
to know that \Alpha\ and \LaTeX\ work well together.  To typeset the 
file you're currently editing in \Alpha, simply choose \acmd{Typeset} 
from the \amenu{Process} submenu or press \key{Cmd T}.  \Alpha\ first 
checks to make sure that any changes to the file have been saved; if 
not, the user is prompted for the appropriate action.  Note that it is 
not necessary to save the document to process the window.  Just click 
the ``No'' button when asked to save the current window, whereupon 
\Alpha\ will pass the contents of the window to the \LaTeX\ 
application and typeset the file automatically.  If the flag 
\acmd{runTeXInBack} is set to true, typesetting will occur in the 
background.

\paragraph{Inverse search --- jumping to a source file}

The inverse operation, switching from \LaTeX\ to \Alpha, depends on 
which \LaTeX\ application you're using.  \OzTeX\ users, for example, 
simply choose the \acmd{Edit} command from \OzTeX's \amenu{Edit} menu 
or press \key{Cmd E} to return to \Alpha.  (By the way, typing `e' in 
response to a \LaTeX\ error message in the \OzTeX\ window throws you 
back into \Alpha\ at the offending line.  The same trick works in 
\CMacTeX\ and \DirectTeX, too.)

Many TeX-related applications allow the possibility of jumping back
to the source file corresponding to a certain situation (an error in 
processing or a mouse click in a dvi window).  This is sometimes
known a ``inverse search''

For example, the xdvi viewer works very well in OS X. If one clicks on
an xdvi window then xdvi attempts to execute a unix command that one
can specify, e,g.  a unix shell script.  Typically, this is used to
open an editor with a specified file and at a specified place.  Thus,
if using Alphatk on MacOS X one can set the environment variable
XEDITOR or EDITOR to
\begin{verbatim}
    /Applications/Alphatk.app/Contents/MacOS/Alphatk +%l %f
\end{verbatim}
or on Windows:
\begin{verbatim}
    C:/Apps/alphatk.exe +%l %f
\end{verbatim}
and clicking on an xdvi window opens the tex source file in alphatk at
the right spot.

If operating with the cygwin environment on Windows (using Xdvi
running in Xwin on top of cygwin), then one should use 
\begin{verbatim}
    C:/Apps/alphatk.exe -cygwin +%l %f
\end{verbatim}
so that Alphatk knows to interpret the path as a cygwin-style path
instead of an ordinary Windows path.

\paragraph{MikTeX}

When using MikTeX, the following should be included in the
``miktex.ini'' file:
\begin{verbatim}
    Editor = C:/Apps/alphatk.exe +"%l" %f
\end{verbatim}

\paragraph{View synchronization}

This is also known as ``forwards search'' or ``synchronicity''.  When
editing any particular .tex file, one can jump to that position in the
typeset document (in a .dvi view, or in \Textures\ view).  This is
known to work with \Textures\ on MacOS, with Yap and Xdvi on Windows.
Normally \key{Control `} is bound to this action.

The commands that Alphatk uses for this, internally, are:

\begin{verbatim}
    yap.exe -1 -s "%l %n" "%P.dvi"
    xdvi -sourceposition "%l %n" "%P.dvi"
\end{verbatim}
where \%l is the linenumber, \%n is the tail of the name of the source
file, and \%P is the name of the dvi file (the tex master file).

\paragraph{More process menu commands}

\begin{commands}
	\acmd{Typeset Clipboard} & \key{Shf Cmd T} \\
	\acmd{Typeset Selection} &
\end{commands}
To typeset the contents of the Clipboard, choose the \acmd{Typeset 
Clipboard} command from the \amenu{Process} submenu or press \key{Shf 
Cmd T}.  This command is handy for typesetting and viewing \TeX\ or 
\LaTeX\ code copied to the Clipboard from other applications such as 
terminal emulators or e-mail clients.

It's also possible to typeset a portion of a document.  Simply select 
(i.e., highlight) the \LaTeX\ code you'd like to typeset and choose 
\acmd{Typeset Selection} from the \amenu{Process} submenu.  \Alpha\ will 
construct a temporary document from the current document's preamble 
and the highlighted text, and pass this virtual document to the \TeX\ 
application to be typeset automatically.

\begin{commands}
	\acmd{dvips file.dvi} & \\
	\acmd{Open file.ps} & \\
	\acmd{View file.ps} & \\
	\acmd{Print file.ps} &
\end{commands}
To convert a \filetype{dvi} file to a \filetype{ps} file, choose the 
\acmd{dvips} command on the \amenu{Process} submenu.  Assuming you 
have the necessary applications installed on your Macintosh, choose 
\acmd{View file.ps} or \acmd{Print file.ps} to view or print the 
resulting \filetype{ps} file.

Once a \filetype{ps} file has been created, you may open a window 
containing the raw PostScript code by choosing \acmd{Open file.ps} on 
the \amenu{Process} submenu.  To see this command, press the \key{Opt} 
key while the \amenu{Process} submenu is down.

\begin{commands}
	\acmd{bibtex file.aux} & \\
	\acmd{Open file.bbl} & \\
	\acmd{makeindex file.idx} & \\
	\acmd{Open file.ind} &
\end{commands}
To run \BibTeX\ or \MakeIndex, choose the corresponding command from 
the \amenu{Process} submenu.  While the \amenu{Process} submenu is 
down, press the \key{Opt} key and choose \acmd{Open file.bbl} or 
\acmd{Open file.ind} to open the file created by \BibTeX\ or 
\MakeIndex, respectively.

\begin{commands}
	\acmd{Open file.log} & \\
	\acmd{Open file.aux} & \\
	\acmd{Open file.toc} & \\
	\acmd{Open file.lof} & \\
	\acmd{Open file.lot} & \\
	\acmd{Open file.idx} & \\
	\acmd{Open file.blg} & \\
	\acmd{Open file.ilg} & \\
	\acmd{Open Any TeX File\ldots} & \key{Shf Cmd O}
\end{commands}
The \amenu{Other Files} submenu on the \amenu{Process} submenu 
provides convenient access to other \LaTeX\ auxiliary files.  Choose 
\acmd{Open Any TeX File} on the \amenu{Other Files} submenu to open 
\emph{any} file in the current directory.

\begin{commands}
	\acmd{Remove Auxiliary Files\ldots} &  \\
	\acmd{Remove Temporary Files} &
\end{commands}
The utility \acmd{Remove Auxiliary Files} interactively removes all 
auxiliary files (\filetype{aux} \filetype{bbl} \filetype{dvi} 
\filetype{glo} \filetype{idx} \filetype{ind} \filetype{lof} 
\filetype{log} \filetype{lot} \filetype{toc} \filetype{blg} 
\filetype{clg} \filetype{ilg} \filetype{ps}) in the current directory.  
Two additional buttons have been added to the dialog: the button 
labeled ``rm ext'' removes all files with the same extension as the 
file displayed in the dialog, and ``rm all'' removes all auxiliary 
files from the current directory without prompting.

\Alpha\ writes all temporary files to \filespec{\$PREFS:tmp:}, which 
makes them easier to remove.  All temporary files are removed once, at 
launch; however, the command \acmd{Remove Temporary Files} on the 
\amenu{Process} submenu removes all temporary files immediately.

\paragraph{Experimental interactive processing}

Not included in the process menu is some new experimental commands
which allow you to execute interactive LaTeX runs from inside
Alpha(tk).

Use \key{Shift Command R} to launch such an interactive session inside
one of Alpha's windows.  If errors occur, you can use `e' to switch
directly to the offending text.

Once this code is robust and tested, and known to work well across
Windows and MacOS X, it will be integrated into the standard process
menu.

\subsubsection{\amenu{Goto} submenu}%
\label{sec:gotoSubmenu}

\begin{commands}
	\acmd{LaTeX} & \key{Shf Cmd S} \\
	\acmd{BibTeX} &  \\
	\acmd{MakeIndex} &
\end{commands}
These commands launch and switch to the corresponding application 
\emph{without} saving and typesetting the current document.  The 
\acmd{LaTeX} command, for instance, is identical to the old 
\acmd{latex} command in \filespec{latex.tcl}~v2.0.

\begin{commands}
	\acmd{Next Template Stop} & \key{Tab} \\
	\acmd{Prev Template Stop} & \key{Shf Tab} 
\end{commands}
As you write your document using the various commands on the \LaTeX\ 
menu, templates are inserted into the text along with tab stops 
(represented by bullets, which may also be inserted with \key{Opt 8}).  
The idea is to type an argument at the current tab stop, press 
\key{Tab} to go to the next tab stop, enter another argument, press 
\key{Tab} again, and so on.  That's what \acmd{Next Tab Stop} and 
\acmd{Prev Tab Stop} do: they move around from tab stop to tab stop.  
Since \acmd{Next Tab Stop} and \acmd{Prev Tab Stop} are bound to the 
\key{Tab} and \key{Shf Tab}, respectively, the menu commands aren't as 
convenient as simply pressing the tab key, but they're included on the 
\LaTeX\ menu for completeness.

NOTE:  Press \key{Opt Tab} to insert a literal tab into the document.

\begin{commands}
	\acmd{Prev Command} & \key{Kpad4} \\
	\acmd{Next Command} & \key{Kpad6} \\
	\acmd{Prev Command Select} & \key{Shf Kpad4} \\
	\acmd{Next Command Select} & \key{Shf Kpad6} \\
	\acmd{Prev Command Select With Args} & \key{Shf Opt Kpad4} \\
	\acmd{Next Command Select With Args} & \key{Shf Opt Kpad6}
\end{commands}
The \acmd{Prev Command} and \acmd{Next Command} commands move the 
cursor to the beginning of the previous or next \LaTeX\ command, while 
\acmd{Prev Command Select} and \acmd{Next Command Select} select the 
previous or next \LaTeX\ command.  Similarly, \acmd{Prev Command 
Select With Args} and \acmd{Next Command Select With Args} select the 
previous or next \LaTeX\ command, along with any command arguments that 
may be present.  Required arguments containing nested braces will not be 
selected, however.  See section~\ref{sec:bugs} for more information 
about this and other \Alpha\ \LaTeX\ bugs.

\begin{commands}
	\acmd{Prev Environment} & \key{Cmd Kpad4} \\
	\acmd{Next Environment} & \key{Cmd Kpad6} \\
	\acmd{Prev Environment Select} & \key{Shf Cmd Kpad4} \\
	\acmd{Next Environment Select} & \key{Shf Cmd Kpad6}
\end{commands}
Like \acmd{Prev Command} and \acmd{Next Command}, these commands 
either move the cursor to the beginning of the previous or next 
\LaTeX\ environment, or select the previous or next \LaTeX\ 
environment.  They are useful for locating or relocating environments.

\begin{commands}
	\acmd{Prev Section} & \key{Cmd Kpad8} \\
	\acmd{Next Section} & \key{Cmd Kpad2} \\
	\acmd{Prev Section Select} & \key{Shf Cmd Kpad8} \\
	\acmd{Next Section Select} & \key{Shf Cmd Kpad2}
\end{commands}
The \acmd{Prev Section} and \acmd{Next Section} commands may be used 
to navigate large files with many sections.  They use the regular 
expression \acmd{funcExprAlt} (which, of course, may be modified) 
discussed in section~\ref{sec:variables}.  The \acmd{Prev Section 
Select} and \acmd{Next Section Select} commands select the previous or 
next section, that is, all the text from one \cmd{section} command to 
the next, and are useful for relocating large blocks of text.

\begin{commands}
	\acmd{Prev Subsection} & \key{Kpad8} \\
	\acmd{Next Subsection} & \key{Kpad2} \\
	\acmd{Prev Subsection Select} & \key{Shf Kpad8} \\
	\acmd{Next Subsection Select} & \key{Shf Kpad2}
\end{commands}
The \acmd{Prev Subsection} and \acmd{Next Subsection} commands are 
similar to \acmd{Prev Section} and \acmd{Next Section} except that 
they also stop at each \cmd{subsection} and \cmd{subsubsection} as 
well.  They use the variable \acmd{funcExpr} discussed in 
section~\ref{sec:variables}.  In \TeX\ mode, these commands take the 
place of \Alpha's generic \acmd{Next Func} and \acmd{Prev Func} 
commands, which are bound to \key{Kpad3} and \key{Kpad1}, 
respectively, in other modes.  Like \acmd{Prev Section Select} and 
\acmd{Next Section Select}, the \acmd{Prev Subsection Select} and 
\acmd{Next Subsection Select} commands select the previous or next 
\cmd{section}, \cmd{subsection}, or \cmd{subsubsection}.

\subsubsection{\amenu{LaTeX Utilities} submenu}%
\label{sec:latexUtilsSubmenu}

\begin{commands}
	\acmd{Choose Command} & \key{Shf Cmd C}
\end{commands}
This command has been removed from \LaTeX\ mode version 4, but will 
hopefully be replaced in the future.

This command provides access to each and every command on the \LaTeX\ 
menu via the keyboard.  It's a multi-step process, where the number of 
steps depend on whether you're using the long or short \LaTeX\ menu: 
first, press \key{Shf Cmd C} and choose a submenu from the list (using 
the arrow keys or by pressing the first letter of a submenu name).  
Next, choose another submenu from the list or the desired command, 
whichever is appropriate.  Continue descending the \LaTeX\ submenus 
until the desired command is found.

\begin{commands}
	\acmd{Delete Tab Stops} & \key{Cmd Tab} \\
	\acmd{Delete Comments} & 
\end{commands}
The \acmd{Delete Tab Stops} command deletes all tab stops (bullets) 
from the current document (or the current selection, if there is one).  
The \acmd{Delete Comments} command deletes all \emph{unnecessary} 
comments from a \LaTeX\ document (which is more difficult than you 
think).  Using the \acmd{Find} dialog, the following three-step manual 
operation (try it!) will remove all comments from the current document:
\begin{center}
	\begin{tabular}{rcc}
		 & search string & replace string  \\
		\hline
		step 1: & \verb*|^[ \t]*%[^\r\n]*\r| & null  \\
		step 2: & \verb*|[ \t]+%[^\r\n]*| & null  \\
		step 3: & \verb*|([^\\](\\\\)*)%[^\r\n]*| & \verb*|\1%|  \\
	\end{tabular}
\end{center}
The utility \acmd{Delete Comments} simply automates this process.
Thanks to Craig Platt
\href{mailto:platt@cc.umanitoba.ca}{platt@cc.umanitoba.ca} for posting
this algorithm in the newsgroup \texttt{comp.text.tex}.

WARNING!  The effects of \acmd{Delete Comments} can not be undone.

\begin{commands}
	\acmd{Convert Quotes} &  \\
	\acmd{Convert Dollar Signs} & 
\end{commands}
If there is a selection, \acmd{Convert Quotes} converts all straight 
quotes to curved quotes (\LaTeX-style) within the selection; 
otherwise, it converts the entire document.

Plain \TeX\ uses dollar signs to delimit math mode and displaymath 
mode.  Since \LaTeX\ inherits most, if not all of plain \TeX's 
functionality, dollar signs work in \LaTeX\ documents, too.  Identical 
left and right delimiters are difficult to parse, however, and so any 
error messages will be misleading at best.  That is why \LaTeX\ has 
its own math mode delimiters and that's why they should be used.  The 
\acmd{Convert Dollar Signs} command replaces all dollar signs in the 
current document (or the current selection, if there is one) with 
appropriate \LaTeX\ syntax.  It does this by making two passes over 
the code, and is therefore somewhat slow on large documents.

\begin{commands}
	\acmd{Compress Text/Math Menus} & 
\end{commands}
The \acmd{Compress Text/Math Menus} commands are toggleable menu items that 
affect the appearance of the \LaTeX\ menu.
the discussion on page~\pageref{pg:compressMenus} for details.


\subsection{Document-related commands}

\subsubsection{\amenu{Documents} submenu}%
\label{sec:documentsSubmenu}

\begin{commands}
	\acmd{New Document} & \key{Shf Cmd N}
\end{commands}
Use this command to open a new window in \TeX\ mode.  Choose \acmd{New 
Document} or press \key{Shf Cmd N} to bring up a dialog with a pop-up 
menu of standard document types.  This will create a new \TeX\ window, 
insert a document of the requested type, and automatically run the 
\acmd{options} command (which is still on the \amenu{Documents} 
submenu).  The old commands \acmd{article}, \acmd{letter}, etc.\ will 
be found on the \amenu{Insert Document} subsubmenu (discussed below).  
Each such command behaves as it did before, that is, it inserts a 
document template into an empty window or wraps the entire contents of 
the current window.

\begin{commands}
	\acmd{article} &  \\
	\acmd{report} &  \\
	\acmd{book} &  \\
	\acmd{letter} &  \\
	\acmd{slides} &  \\
	\acmd{generic\ldots} & 
\end{commands}
Choosing one of these document templates from the \amenu{Documents} 
submenu either inserts the desired template at the insertion point, or 
if there is a current selection, the selection is wrapped up inside 
the chosen template.  In either case, the insertion point is 
positioned at the beginning of the template where the user may enter 
any specific document class options that may be required (standard 
options include \texttt{11pt}, \texttt{twoside} and 
\texttt{twocolumn}, for example).  If none are desired, simply skip 
over this part of the template (it's okay to leave the square brackets 
empty).  See the \acmd{options} command below for more information on 
document class options.

\begin{commands}
	\acmd{options\ldots} &  \\
	\acmd{usepackage} & \key{Ctl Opt U}
\end{commands}
The \acmd{options} command presents the user with a dialog box and a 
list of standard document class options.  Choosing one of these 
options or typing a name into the text box of the dialog inserts the 
chosen option into the current document at the appropriate place.  See 
the bug list in section~\ref{sec:bugs} for a caveat, however.

When you insert a \LaTeX\ document template into the current window, 
you get one \cmd{usepackage} command by default.  To insert another 
\cmd{usepackage} immediately after the \cmd{documentclass} command, 
choose the \acmd{usepackage} command on the \amenu{Documents} submenu 
or press \key{Ctl Opt U}.

\begin{commands}
	\acmd{filecontents\ldots} &  \\
	\acmd{filecontents All} & 
\end{commands}
To facilitate file transfer, \LaTeXe\ now has a \env{filecontents} 
environment that contains the source of a \LaTeX\ auxiliary file or 
input file.  Issuing this command brings up a standard file dialog.  
After locating the file to be included, \filespec{latex.tcl} wraps the 
file inside a \env{filecontents} environment and inserts it at the 
beginning of the document.

There is also a \acmd{filecontents All} command that scans the 
current document and prepends one \env{filecontents} environment for 
each custom package or class file in the current folder.  Local files 
read by \cmd{input} or \cmd{include} are also attached, as well as 
\filetype{bib} and \filetype{bst} files.

\begin{commands}
	\acmd{Rebuild Documents Submenu} & 
\end{commands}
This command rebuilds the \amenu{Documents} submenu on-the-fly.  It's 
a temporary fix until I think of a better way to handle the 
\amenu{Packages} submenu.  The \amenu{Packages} submenu contains a 
list of all packages known to the \TeX\ application.  Choosing one of 
these packages inserts the corresponding \cmd{usepackage} command into 
the preamble of the current document.  To build this submenu, enable 
the flag \acmd{buildPkgsSubmenu} as described in 
section~\ref{sec:flags}, and then choose \acmd{Rebuild Documents 
Submenu} on the \amenu{Documents} submenu.

\subsubsection{\amenu{Page Layout} submenu}

\begin{commands}
	\acmd{maketitle} & 
\end{commands}
\LaTeX's \cmd{maketitle} command formats a title page with information 
provided by the user.  Choosing this command from the \LaTeX\ menu 
inserts a title page template into the current document just after the 
\cmdarg{begin}{document} command.

\begin{commands}
	\acmd{abstract} &  \\
	\acmd{titlepage} & 
\end{commands}
The \env{abstract} and \env{titlepage} environments contain the text 
of an abstract and title page, respectively.  The latter differs from 
\cmd{maketitle} in that the user is totally responsible for the format 
of the title page.

\begin{commands}
	\acmd{pagestyle\ldots} &  \\
	\acmd{thispagestyle\ldots} &  \\
	\acmd{pagenumbering\ldots} & 
\end{commands}
The \acmd{pagestyle} and \acmd{thispagestyle} commands control what 
appears in the header and footer of the current document.  The user is 
presented with a list of standard formats from which to choose.  The 
\acmd{pagenumbering} command is for choosing the style of the page 
numbers, and is also interactive.

\begin{commands}
	\acmd{twocolumn} &  \\
	\acmd{onecolumn} & 
\end{commands}
These are simple declarations that tell \LaTeX\ to begin formatting 
the output in two or one column format, respectively.

\subsubsection{\amenu{Sectioning} submenu}

\begin{commands}
	\acmd{part} &  \\
	\acmd{chapter} &  \\
	\acmd{section} &  \\
	\acmd{subsection} &  \\
	\acmd{subsubsection} &  \\
	\acmd{paragraph} &  \\
	\acmd{subparagraph} & 
\end{commands}
All \LaTeX\ sectioning commands are available from the 
\amenu{Sectioning} submenu, the most common commands being the 
\acmd{chapter}, \acmd{section}, and \acmd{subsection} commands.  The 
corresponding \LaTeX\ command is inserted at the insertion point.  The 
current selection, if there is one, is assumed to be the name of the 
section and wrapped up inside curly braces.  The resulting declaration 
is \emph{not} automatically followed by a carriage return since the 
user has the option of putting a label (or whatever) on the same line.

\begin{commands}
	\acmd{appendix} & 
\end{commands}
Unlike the other sectioning commands, this command does not have an 
argument.  It simply tells \LaTeX\ to start numbering differently.  
The \cmd{appendix} declaration only makes sense in the context of a 
long document such as a book.


\subsection{Paragraph mode commands}

\subsubsection{\amenu{Text Style} submenu}

The following text style commands each take an argument, namely, the 
text to be formatted in the given style.  For large amounts of text, 
use the corresponding declarations listed on p.~37 of the \LaTeX\ 
book.

\begin{commands}
	\acmd{emph} & \key{Ctl Opt E} \\
	\acmd{underline} & \key{Ctl Cmd U}
\end{commands}
Short for ``emphasized'', the \acmd{emph} command is perhaps the most 
often used \LaTeX\ text style.  If the surrounding text has the 
upright shape (see below), then \LaTeX\ typesets emphasized text in 
italics.  If the surrounding text is italicized, then emphasized text 
will be upright.  The so-called ``italic correction'' is handled 
automatically by this command.

Although underlined text is not used much anymore, the corresponding 
command is included here for completeness.  The \acmd{underline} 
command may also be used in math mode and therefore also appears on 
the \amenu{Grouping} submenu.  See section~\ref{sec:groupingSubmenu}.

\begin{commands}
	\acmd{textup} &  \\
	\acmd{textit} & \key{Ctl Opt I} \\
	\acmd{textsl} & \key{Ctl Opt S} \\
	\acmd{textsc} & \key{Ctl Opt H}
\end{commands}
These four commands specify the \emph{shape} of their respective 
arguments.  They call for upright text, italics, slanted text, and 
small caps, respectively.  Upright is the default.

\begin{commands}
	\acmd{textmd} &  \\
	\acmd{textbf} & \key{Ctl Opt B}
\end{commands}
These commands specify an attribute called the \emph{series} of the 
corresponding font.  They call for medium and boldfaced text, 
respectively.  Medium is the default.

\begin{commands}
	\acmd{textrm} & \key{Ctl Opt R} \\
	\acmd{textsf} & \key{Ctl Opt W} \\
	\acmd{texttt} & \key{Ctl Opt Y}
\end{commands}
The third and final component of any given font is the \emph{family}.  
There are three families: roman, sans serif, and typewriter.  Roman is 
the default.

\begin{commands}
	\acmd{textnormal} & 
\end{commands}
Regardless of the surrounding text, the argument of \cmd{textnormal} 
is typeset in the default style, that is, upright, medium, and roman.

\begin{commands}
	\acmd{em} & \key{Shf Ctl Opt E} \\
	\acmd{upshape} &  \\
	\acmd{itshape} & \key{Shf Ctl Opt I} \\
	\acmd{slshape} & \key{Shf Ctl Opt S} \\
	\acmd{scshape} & \key{Shf Ctl Opt H} \\
	\acmd{mdseries} &  \\
	\acmd{bfseries} & \key{Shf Ctl Opt B} \\
	\acmd{rmfamily} & \key{Shf Ctl Opt R} \\
	\acmd{sffamily} & \key{Shf Ctl Opt W} \\
	\acmd{ttfamily} & \key{Shf Ctl Opt Y} \\
	\acmd{normalfont} & 
\end{commands}
These commands are the declarative counterparts of the previously 
mentioned text style commands.  Typically, they are used for large 
chunks of text, say, entire paragraphs.  (Note: the declarative 
versions do not apply an italic correction.  See the \LaTeX\ manual 
for usage and examples.)  To access these commands, press the \key{Shf} 
key with the \amenu{Text Style} submenu down.

\subsubsection{\amenu{Text Size} submenu}

\begin{commands}
	\acmd{tiny} & \key{Ctl Opt 1} \\
	\acmd{scriptsize} & \key{Ctl Opt 2} \\
	\acmd{footnotesize} & \key{Ctl Opt 3} \\
	\acmd{small} & \key{Ctl Opt 4} \\
	\acmd{normalsize} & \key{Ctl Opt 5} \\
	\acmd{large} & \key{Ctl Opt 6} \\
	\acmd{Large} & \key{Ctl Opt 7} \\
	\acmd{LARGE} & \key{Ctl Opt 8} \\
	\acmd{huge} & \key{Ctl Opt 9} \\
	\acmd{Huge} & \key{Ctl Opt 0}
\end{commands}
These commands declare the text font size.  They affect the entire 
document unless surrounded by braces, so the menu commands 
automatically insert braces.  If you want the entire document set in a 
certain font size, insert a class option with the \acmd{options} 
command (see section~\ref{sec:documentsSubmenu}).

\subsubsection{\amenu{International} submenu}

\filespec{latex.tcl} implements about half of \LaTeX's full palette of 
international symbols and accents (if you can think of ways to get the 
rest of these on the \LaTeX\ menu, please let me know!).  See 
Tables~3.1 and 3.2 on pp.~38--39 of the \LaTeX\ book for a complete 
list.

\subsubsection{\amenu{Environments} submenu}%
\label{sec:envSubmenu}

\begin{commands}
	\acmd{itemize\ldots} & \key{Opt F7} \\
	\acmd{enumerate\ldots} & \key{Shf Opt F7} \\
	\acmd{description\ldots} & \key{Ctl Opt F7} \\
	\acmd{thebibliography\ldots} &  \\
	\acmd{slide} & \key{Opt F8} \\
	\acmd{overlay} & \key{Shf Opt F8} \\
	\acmd{note} & \key{Ctl Opt F8} \\
	\acmd{figure} & \key{Opt F9} \\
	\acmd{table} & \key{Shf Opt F9} \\
	\acmd{tabular\ldots} & \key{Ctl Opt F9} \\
	\acmd{verbatim} & \key{Opt F10} \\
	\acmd{quote} & \key{Shf Opt F10} \\
	\acmd{quotation} & \key{Ctl Opt F10} \\
	\acmd{verse} &  \\
	\acmd{center} & \key{Opt F11} \\
	\acmd{flushleft} & \key{Shf Opt F11} \\
	\acmd{flushright} & \key{Ctl Opt F11} \\
	\acmd{general\ldots} & \key{Opt F12}
\end{commands}
One of the most useful of \filespec{latex.tcl}'s many features is its 
ability to insert skeletal templates for multi-line environments (that 
is, \LaTeX\ constructs delimited by a \cmd{begin}\ldots\cmd{end} pair).  
These may be inserted anywhere in the document (even in the middle of 
a line), complete with tab stops and appropriate indentation.  In some 
cases (like \acmd{itemize}), the user is asked to specify the number 
of rows desired, after which the program generates the corresponding 
environment body complete with indentation and tab stops.  Some 
environment commands (like \acmd{tabular}) also prompt the user for 
the desired number of columns.  There's even a \acmd{general} command 
for inserting user-defined environments on-the-fly.

The \acmd{figure} command deserves special mention.  Choosing this 
command from the \amenu{Environments} submenu (or by pressing \key{Opt 
F9}) brings up a dialog with a pop-up menu of box-making macros, one 
for every macro name stored in the \TeX\ mode variable 
\acmd{boxMacroNames} (see section~\ref{sec:variables}).  With the 
mouse or arrow keys, choose one of these macro names and click ``OK'' 
to insert the corresponding \env{figure} environment at the insertion 
point, or leave the text box blank to wrap a \env{figure} environment 
around the current selection (if there is one).  If only one macro 
name is stored in \acmd{boxMacroNames}, the dialog is automatically
circumvented and the \env{figure} environment is inserted at the 
insertion point without prompting.

Note: The \amenu{Environments} submenu seeks to mimic the 
corresponding command keys.  Each group of environments on this 
submenu has been assigned a different function key, beginning with 
\key{F7}.  The \acmd{general} environment, for instance, is bound to 
\key{Opt F12}.  See section~\ref{sec:keys} for more information.

\subsubsection{\amenu{Boxes} submenu}

\begin{commands}
	\acmd{mbox} & \key{Ctl Opt M} \\
	\acmd{makebox} &  \\
	\acmd{fbox} &  \\
	\acmd{framebox} & 
\end{commands}
Perhaps the most useful box-making command is \cmd{mbox}, which 
formats its argument in LR mode, a restricted form of paragraph mode 
impervious to line breaks.  The \cmd{mbox} command is especially 
useful for inserting a bit of plain text in the middle of a math 
formula (see the \textsl{\LaTeX} book for examples).  The \cmd{makebox} 
command is a generalized form of \cmd{mbox}, which takes the width and 
height of the box as additional arguments.

The commands \cmd{fbox} and \cmd{framebox} are analogous to \cmd{mbox} 
and \cmd{makebox} except that a rectangular frame is drawn around the 
box.

\begin{commands}
	\acmd{newsavebox} &  \\
	\acmd{sbox} &  \\
	\acmd{savebox} &  \\
	\acmd{usebox} & 
\end{commands}
A ``savebox'' is a bin for storing text, graphics, formulas, or 
whatever.  The argument to \cmd{sbox} or \cmd{savebox} is typeset 
\emph{once} and may be recalled later, any number of times, via 
\cmd{usebox}.

\begin{commands}
	\acmd{raisebox} & 
\end{commands}
This box-making command takes a vertical offset as one of its arguments.

\begin{commands}
	\acmd{parbox} &  \\
	\acmd{minipage} & 
\end{commands}
The primary argument of \LaTeX's \cmd{parbox} command or 
\env{minipage} environment is typeset in paragraph mode.  \cmd{parbox} 
is for small amounts of text, while the \env{minipage} environment is 
for large blocks of text.

\begin{commands}
	\acmd{rule} & 
\end{commands}
The \cmd{rule} command makes a box filled with ink.  For example,
\begin{verbatim}
  \newcommand{\filledsquare}{\rule[0.125ex]{1.3ex}{1.3ex}}
\end{verbatim}
makes a black square approximately the same size as \LaTeX's open \cmd{Box}.  
(There is an analogous command called \cmd{blacksquare} defined in the AMS 
symbol package.)

\subsubsection{\amenu{Miscellaneous} submenu}%
\label{sec:miscSubmenu}

\begin{commands}
	\acmd{verb} & \key{Ctl Opt V} \\
	\acmd{footnote} & \key{Ctl Opt F} \\
	\acmd{marginal note} & \key{Ctl Opt N}
\end{commands}
All the above commands wrap the current selection, if there is one.

\begin{commands}
	\acmd{label} & \key{Ctl Opt L} \\
	\acmd{ref} & \key{Ctl Opt X} \\
	\acmd{pageref} & \key{Ctl Opt P} \\
	\acmd{cite} & \key{Ctl Opt C} \\
	\acmd{nocite} & \key{Shf Ctl Opt C}
\end{commands}
These commands do more than simply insert the corresponding \LaTeX\ 
command.  For instance, press \key{Ctl Opt X} or \key{Ctl Opt P} to 
insert a \cmd{ref} or \cmd{pageref} command, respectively.  The 
inserted command will contain the argument of the nearest \cmd{label} 
command.  Continue pressing \key{Ctl Opt X} or \key{Ctl Opt P} to cycle 
through all the \cmd{label} commands in your document.

\begin{commands}
	\acmd{item} & \key{Ctl Opt J}
\end{commands}
Simply press \key{Ctl Opt J} inside an \env{itemize}, \env{enumerate}, 
\env{description}, or \env{thebibliography} environment to insert an 
item of the appropriate type at the insertion point.

\begin{commands}
	\acmd{quotes} & \key{Ctl Opt \straightquote} \\
	\acmd{double quotes} & \key{Shf Ctl Opt \straightquote}
\end{commands}
The \filespec{latex.tcl} macro package incorporates a ``smart quotes'' 
feature originally implemented by an unknown author (see the code in 
\filespec{latexSmart.tcl}) that makes the typing of quoted material 
totally transparent.  Just use the quote key as you would for plain 
text files.  Consequently, the \acmd{quotes} and \acmd{double quotes} 
commands are primarily used for quoting existing text.

\begin{commands}
	\acmd{ellipsis} &  \\
	\acmd{en-dash} &  \\
	\acmd{em-dash} &  \\
	\acmd{TeX logo} &  \\
	\acmd{LaTeX logo} &  \\
	\acmd{Latex2e logo} &  \\
	\acmd{date} & 
\end{commands}
These are a few of the text-related \LaTeX\ commands that I've found 
useful from time to time.

\begin{commands}
	\acmd{dag} &  \\
	\acmd{ddag} &  \\
	\acmd{section mark} &  \\
	\acmd{paragraph mark} &  \\
	\acmd{copyright} &  \\
	\acmd{pounds} & 
\end{commands}
The previous six commands may be used in any mode, including math mode.


\subsection{Math mode commands}

\subsubsection{\amenu{Math Mode} submenu}%
\label{sec:mathModeSubmenu}

\begin{commands}
	\acmd{TeX math} &  \\
	\acmd{TeX displaymath} &  \\
	\acmd{LaTeX math} & \key{Ctl Cmd M} or \key{Ctl Cmd 4} \\
	\acmd{LaTeX displaymath} & \key{Ctl Opt Cmd M} or \key{Ctl Opt Cmd 4}
\end{commands}
Math mode may be invoked in a number of ways.  Many \TeX nical typists 
rely exclusively on \TeX's use of dollar signs and almost always key 
in their documents horizontally from left to right.  Others have 
adopted \LaTeX's tendency to prefer vertical constructions 
(environments).  Still others have settled on some combination of 
these, using whichever seems comfortable or convenient at the time.  
Whatever your approach to mathematical typesetting, there's something 
for everybody in \filespec{latex.tcl}, designed to simplify the input 
of complex mathematical formulas.

Four math modes are available for normal, left-to-right input.  These 
are called \acmd{TeX math} \verb|$|\ldots\verb|$| and \acmd{TeX 
displaymath} \verb|$$|\ldots\verb|$$|, along with their 
corresponding \LaTeX\ equivalents called \acmd{LaTeX math} 
\cmd{(}\ldots\cmd{)} and \acmd{LaTeX displaymath} 
\cmd{[}\ldots\cmd{]}.  The \LaTeX\ versions are logically equivalent 
to the multi-line \env{math} and \env{displaymath} environments (see 
below).  The latter have the advantage that 1) they are often more 
readable in source form, and 2) they are more easily changed (by 
simply replacing keywords) as the document evolves.

NOTE: The above command keys automatically switch from 
\acmd{LaTeX math} and \acmd{LaTeX displaymath} to \acmd{TeX math} and 
\acmd{TeX displaymath}, respectively, when the flag 
\acmd{useDollarSigns} is set to true.

\textbf{Tip:} Get into the habit of pressing \key{Ctl Cmd M} or 
\key{Ctl Opt Cmd M} when composing in-line equations, since there is 
less chance of inadvertantly omitting a dollar sign if you do.

\subsubsection{\amenu{Math Style} submenu}

\begin{commands}
	\acmd{mathit} & \key{Ctl Opt Cmd I} \\
	\acmd{mathrm} & \key{Ctl Opt Cmd R} \\
	\acmd{mathbf} & \key{Ctl Opt Cmd B} \\
	\acmd{mathsf} & \key{Ctl Opt Cmd W} \\
	\acmd{mathtt} & \key{Ctl Opt Cmd Y} \\
	\acmd{mathcal} & \key{Ctl Opt Cmd C} \\
	\acmd{displaystyle} & \key{Ctl Opt Cmd D} \\
	\acmd{textstyle} & \key{Ctl Opt Cmd T} \\
	\acmd{scriptstyle} & \key{Ctl Opt Cmd S} \\
	\acmd{scriptscriptstyle} & 
\end{commands}
The next submenu on the \LaTeX\ menu is called \amenu{Math Style}, 
with commands for math italic, roman, boldface, sans serif, typewriter 
and calligraphic typefaces, as well as declarations for 
\cmd{displaystyle}, \cmd{textstyle}, \cmd{scriptstyle}, and 
\cmd{scriptscriptstyle}.  The latter command quartet are sometimes 
needed to override \LaTeX's default math style.  (The \env{array} 
environment, for example, insists on enabling \cmd{textstyle} 
regardless of the surrounding environment.)

\subsubsection{\amenu{Math Environments} submenu}

\begin{commands}
	\acmd{math} & \key{Opt F5} \\
	\acmd{displaymath} & \key{Shf Opt F5} \\
	\acmd{equation} & \key{Ctl Opt F5} \\
	\acmd{eqnarray\texttt{*}} & \key{Shf Opt F6} \\
	\acmd{eqnarray} & \key{Ctl Opt F6} \\
	\acmd{array} & \key{Opt F6} \\
	\acmd{general} & \key{Opt F12}
\end{commands}
Besides the \env{math} and \env{displaymath} environments discussed in 
section~\ref{sec:mathModeSubmenu}, other multi-line math environments 
(\env{equation}, \env{array}, \env{eqnarray}, and \env{eqnarray*}) are 
also available.  Each is mutually exclusive (that is, one may not be 
nested inside the other) except for the \env{array} environment which 
\emph{must} be nested inside some other math environment.  (It took me 
a long time to come to grips with this apparent anomaly).  There's 
also a \acmd{general} environment command, which is exactly the same 
command found on the \amenu{Text Style} submenu.

\subsubsection{\amenu{Formulas} submenu}%
\label{sec:formulasSubmenu}

\begin{commands}
	\acmd{subscript} & \key{\_} (if \acmd{smartScripts} is true) \\
	\acmd{superscript} & \key{\^{}} (if \acmd{smartScripts} is true) \\
	\acmd{frac} & \key{Ctl Cmd F} \\
	\acmd{sqrt} & \key{Ctl Cmd R} \\
	\acmd{nth root} &  \\
	\acmd{one parameter} & \key{Ctl Cmd 1} \\
	\acmd{two parameters} & \key{Ctl Cmd 2}
\end{commands}
The \amenu{Formulas} submenu contains \LaTeX\ commands commonly used 
to build up even the simplest mathematical expressions.  There are 
commands for typesetting subscripts and superscripts, fractions (which 
used to be difficult to typeset), square roots, and arbitrary $n$th 
roots.  There are also one and two-parameter \LaTeX\ commands, which 
allow the user to type in a command name on-the-fly.  Next to 
\filespec{latex.tcl}'s environment commands, the formula commands are 
most useful.  (In fact, it pays to memorize their command key 
equivalents.)

While we're talking about the \amenu{Formulas} submenu, let me say a 
little bit about \filespec{latex.tcl}'s ability to parse fractions.  
How many times have you found yourself wanting to recast a 
horizontally typeset fraction such as 
\begin{verbatim}
  $x = (-b \pm \sqrt{b^2 - 4ac})/(2a)$
\end{verbatim}
in a corresponding ``vertical'' form 
\begin{verbatim}
  $$x = \frac{(-b \pm \sqrt{b^2 - 4ac})}{(2a)}$$
\end{verbatim}
Obviously, such an operation involves a lot of cutting and pasting, 
and I used to avoid it like the plague.  Well, now all you have to do 
is select the text you want converted (in this case, all the text 
inside the dollar signs except ``\texttt{x = }``) and then choose the 
\acmd{frac} command from the \amenu{Formulas} submenu on \LaTeX\ menu.  
The rest is automatic.  (Now if only I could get it to automatically 
remove the redundant parentheses\ldots)

\subsubsection{\amenu{Greek} submenu}

One of the longest of \filespec{latex.tcl}'s submenus contains the 
entire Greek alphabet, including both lower and upper-case letters 
(hold down the \key{Opt} key while the \amenu{Greek} submenu is down 
to see the latter), plus a handful of lower-case ``italicized'' 
letters (\cmd{varepsilon}, \cmd{vartheta}, \cmd{varpi}, \cmd{varrho}, 
\cmd{varsigma}, and \cmd{varphi}).  To type a \amenu{Greek} command at 
the keyboard, press \key{Ctl M} \key{letter}, where \key{letter} is 
the same key assigned to that letter by the Macintosh Symbol font.  See 
the file \filespec{latex\_bindings.tex} for a useful summary.

NOTE: There are two \amenu{Greek} submenus.  While one is down, press 
a modifier key (such as \key{Opt}) to see the alternate menu.

\subsubsection{\amenu{Binary Operators} and \amenu{Relations} submenus}

Plain \TeX\ defines an incredible variety of mathematical symbols, 
each transparently available to the \LaTeX\ user.  All of these 
symbols have been implemented in this version of \filespec{latex.tcl}.

There are two \amenu{Relations} menus.  While one is down, press 
a modifier key (such as \key{Opt}) to see the alternate menu.

\subsubsection{\amenu{Arrows}, \amenu{Dots}, and \amenu{Symbols} submenus}

A quick glance at the \LaTeX\ book shows a wide assortment of arrows, 
dots, and miscellaneous mathematical symbols.  Starting with v2.2, all 
of these have been implemented in \filespec{latex.tcl}.  See the 
\amenu{Arrows}, \amenu{Dots}, and \amenu{Symbols} submenus for 
exhaustive lists of available commands.

NOTE: There are two \amenu{Arrows} menus.  While one is down, press a 
modifier key (such as \key{Opt}) to see the alternate menu.

\subsubsection{\amenu{Functions} submenu}

All of \TeX's so-called ``log-like'' functions (\cmd{exp} and 
\cmd{sin}, for instance) have been implemented in this version of 
\filespec{latex.tcl}.  Some of these commands (\acmd{lim}, \acmd{inf}, 
\acmd{sup}, \acmd{liminf}, \acmd{limsup}, \acmd{max}, and \acmd{min}) 
automatically insert a subscript.  Only \acmd{lim} has a command key, 
namely, \key{Ctl Cmd L}.

\subsubsection{\amenu{Large Operators} submenu}

\begin{commands}
	\acmd{sum} & \key{Ctl Cmd S} \\
	\acmd{prod} & \key{Ctl Cmd P} \\
	\acmd{coprod} &  \\
	\acmd{int} & \key{Ctl Cmd I} \\
	\acmd{oint} &  \\
	\acmd{bigcup} &  \\
	\acmd{bigcap} &  \\
	\acmd{bigsqcup} &  \\
	\acmd{bigvee} &  \\
	\acmd{bigwedge} &  \\
	\acmd{bigodot} &  \\
	\acmd{bigotimes} &  \\
	\acmd{bigoplus} &  \\
	\acmd{biguplus} & 
\end{commands}
The \filespec{latex.tcl} macro package also provides support for 
\TeX's so-called ``large operators''.  Commands such as \acmd{sum} 
\key{Ctl Cmd S}, \acmd{prod} \key{Ctl Cmd P}, \acmd{int} \key{Ctl Cmd 
I}, \acmd{bigcup}, \acmd{bigcap}, \acmd{bigvee}, and \acmd{bigwedge} 
may be found on the \amenu{Large Operators} submenu.

\subsubsection{\amenu{Delimiters} submenu}

\begin{commands}
	\acmd{parentheses} &  \\
	\acmd{brackets} &  \\
	\acmd{braces} &  \\
	\acmd{vertical bars} &  \\
	\acmd{other delims\ldots} &  \\
	\acmd{half-open interval} &  \\
	\acmd{half-closed interval} &  \\
	\acmd{big parentheses} &  \\
	\acmd{big brackets} &  \\
	\acmd{big braces} &  \\
	\acmd{big vertical bars} &  \\
	\acmd{other big delims\ldots} &  \\
	\acmd{big left brace} &  \\
	\acmd{other mixed big delims\ldots} & 
\end{commands}
\TeX\ is particularly adept at ``delimiting'' arbitrary-sized 
mathematical expressions.  Examples include parenthesized equations, 
matrices, and determinants.  Since the left and right delimiters need 
not be of the same type, there are a host of options from which to 
choose, which presents an interesting design problem.  A workable 
compromise was achieved by implementing a handful of common delimiters 
explicitly, and then providing access to other more esoteric 
combinations via dialogs.  Consequently, commands for \acmd{big 
parentheses}, \acmd{big brackets}, \acmd{big braces}, and \acmd{big 
vertical bars} (i.e., absolute value signs) will be found on the 
\amenu{Delimiters} submenu, along with a \acmd{big left brace} 
(commonly used to define multi-part functions or systems of 
equations), as well as commands called \acmd{other big delims} and 
\acmd{other mixed big delims}.  The latter two commands are 
interactive---the user either types the delimiter name directly into a 
text box or chooses the desired name from a pop-up menu of available 
options.  Also on the \amenu{Delimiters} submenu are normal-sized 
\acmd{parentheses}, \acmd{brackets}, \acmd{braces}, \acmd{vertical 
bars}, and other fixed-size delimiters.

\begin{commands}
	\acmd{multi-line big parentheses} &  \\
	\acmd{multi-line big brackets} &  \\
	\acmd{multi-line big braces} &  \\
	\acmd{multi-line big vertical bars} &  \\
	\acmd{other multi-line big delims\ldots} &  \\
	\acmd{multi-line big left brace} &  \\
	\acmd{other multi-line mixed big delims\ldots} & 
\end{commands}
All of the big delimiters have multi-line counterparts (i.e, a 
vertical, as opposed to a horizontal construct).  To access these 
commands, press the \key{Opt} key while the \amenu{Delimiters} submenu 
is down.

\subsubsection{\amenu{Math Accents} submenu}

\begin{commands}
	\acmd{acute} & \key{Ctl Cmd A} \\
	\acmd{bar} & \key{Ctl Cmd B} \\
	\acmd{breve} &  \\
	\acmd{check} & \key{Ctl Cmd C} \\
	\acmd{dot} & \key{Ctl Cmd D} \\
	\acmd{ddot} &  \\
	\acmd{grave} & \key{Ctl Cmd G} \\
	\acmd{hat} & \key{Ctl Cmd H} \\
	\acmd{tilde} & \key{Ctl Cmd T} \\
	\acmd{vec} & \key{Ctl Cmd V} \\
	\acmd{widehat} &  \\
	\acmd{widetilde} &  \\
	\acmd{imath} &  \\
	\acmd{jmath} & 
\end{commands}
Math accents (not to be confused with diacritical marks used in 
paragraph mode) are accessed from a submenu of the same name.  There 
are commands for hats, bars, tildes, vectors, dots, etc., plus wide 
hats and tildes, most of which have command keys.  There are also 
commands for dotless versions of the letters ``$i$'' and ``$j$'' used 
in conjunction with these accents.  Insofar as possible, the macros 
check to make sure that only single characters are being accented, or 
in the case of wide accents, three or fewer characters.

\subsubsection{\amenu{Grouping} submenu}%
\label{sec:groupingSubmenu}

\begin{commands}
	\acmd{underline} & \key{Ctl Cmd U} \\
	\acmd{overline} & \key{Ctl Cmd O} \\
	\acmd{underbrace} & \key{Ctl Opt Cmd U} \\
	\acmd{overbrace} & \key{Ctl Opt Cmd O} \\
	\acmd{overrightarrow} &  \\
	\acmd{overleftarrow} &  \\
	\acmd{stackrel} & 
\end{commands}
The \amenu{Grouping} submenu has commands for underlining and 
overlining, and related commands that produce underbraces and 
overbraces.  There's also a command called \acmd{stackrel} used to 
construct compound operators via vertical stacking (see p.~50 of the 
\LaTeX\ book for more information).

\subsubsection{\amenu{Spacing} submenu}

\begin{commands}
	\acmd{neg thin} &  \\
	\acmd{thin} &  \\
	\acmd{medium} &  \\
	\acmd{thick} &  \\
	\acmd{quad} &  \\
	\acmd{qquad} &  \\
	\acmd{hspace} &  \\
	\acmd{vspace} &  \\
	\acmd{hfill} &  \\
	\acmd{vfill} &  \\
	\acmd{smallskip} &  \\
	\acmd{medskip} &  \\
	\acmd{bigskip} & 
\end{commands}
The \amenu{Spacing} submenu provides for various types of horizontal 
and vertical spacing.  There are commands for negative thin, thin, 
medium, and thick amounts of whitespace, and additional commands for 
inserting the traditional typesetter's quad (1em) and double quad.  
Arbitrary horizontal whitespace, defined via \LaTeX's \cmd{hspace} 
command, and vertical whitespace via \cmd{vspace}, may also be 
inserted from the \amenu{Spacing} submenu.  \LaTeX's ``fill'' commands 
will also be found on this submenu, as well as \cmd{smallskip}, 
\cmd{medskip}, and \cmd{bigskip}.

\section{The \amenu{funcs} menu}
\label{sec:funcs}

The pop-up menu activated by pressing the ``\{\}'' icon on the tool 
bar at the right of each window is called the \amenu{funcs} menu.  In 
\TeX\ mode (and other modes as well), the \amenu{funcs} menu gives an 
outline of the current document and provides a way to quickly navigate 
a long file.  Simply press the ``\{\}'' icon to build the 
\amenu{funcs} pop-up menu on-the-fly.  The current document will be 
scanned and the titles of all sections and subsections will be placed 
on the menu.

\section{The \amenu{mark} menu}
\label{sec:mark}

The pop-up menu activated by pressing the ``\texttt{M}'' icon on the 
tool bar at the right of each window is called the \amenu{mark} menu.  
In \TeX\ mode (and other modes as well), the mark menu gives an 
outline of the current document and provides a way to quickly navigate 
a long file.  Simply press the ``\texttt{M}'' icon and choose the 
\acmd{Mark File} command.  The contents of the current document will 
be scanned and the titles of all chapters, sections, and subsections 
will be placed on the menu.  Files that are \cmd{include}'d or 
\cmd{input}'ed will also appear on the \amenu{mark} menu.  Note that 
the mark menu is static, that is, if you change the structure of the 
current document, you must choose the \acmd{Mark File} command 
again.

Note: Currently, there is considerable overlap between the 
\amenu{funcs} menu and the \amenu{mark} menu in terms of 
functionality.  This will change in future versions of \Alpha.


\chapter{Command Keys}

Menus are great at first, but eventually the tendency is to move away 
from menus towards commands keys.  This can significantly speed the 
input process.  Few of us, however, are inclined to memorize more than 
a couple dozen such keystrokes unless continually prompted with 
reminders.  Thus most command key equivalents are displayed on the 
\LaTeX\ menu in full view.

\section{Tips}
\label{sec:keys}

A few remarks will help you remember the many command key sequences.  
All paragraph mode commands (\cmd{textbf}, \cmd{footnote}, etc.)  
begin with \key{Ctl Opt} or \key{Shf Ctl Opt}, whereas all math mode 
commands begin with \key{Ctl Cmd} or \key{Ctl Opt Cmd}.  All 
environments are bound to some modified function key.  Sometimes the 
\key{Shf} key reverses the orientation of an existing key (as in the 
case of \acmd{Next Tab Stop} and \acmd{Prev Tab Stop} described in 
section~\ref{sec:gotoSubmenu}) or acts as a selection key (see, for 
example, \acmd{Next Environment} and \acmd{Next Environment Select} in 
section~\ref{sec:gotoSubmenu}).  Knowing these simple facts helps 
tremendously.

See section~\ref{sec:docs} for pointers to useful command key summaries.

\section{Double-clicking}
\label{sec:dblClicking}

The arguments of certain \LaTeX\ commands are 
command-double-clickable, that is, you hold down the \key{Cmd} key 
while double-clicking the argument of certain \LaTeX\ commands.  
These commands are underlined and therefore easily recognized in 
your document.  When you command-double-click the required argument 
of a \cmd{ref} command, for example, the cursor jumps to the 
corresponding \cmd{label}.  (Note: Press \key{Ctl .} to return to 
the original cursor position.)  Similarly, when you cmd-dbl-click 
the required argument of a \cmd{cite} command, the cursor jumps to 
the corresponding \cmd{bibitem} if the document contains a 
\env{thebibliography} environment; otherwise, the arguments of a 
\cmd{bibliography} command are sequentially searched until the 
\filetype{bib} file containing the target item is found.  This 
\filetype{bib} file is then opened and the cursor jumps to the 
target item.

\textbf{Tip:} By default, \cmd{ref}, \cmd{pageref}, \cmd{cite}, and 
\cmd{nocite} commands are cmd-dbl-clickable.  If you use a package 
that defines other \cmd{ref}-like or \cmd{cite}-like commands, modify 
the \TeX\ mode variables \acmd{refCommands} or \acmd{citeCommands} 
discussed in section~\ref{sec:variables}.

Other \LaTeX\ commands may also be command-double-clicked, for 
example, \cmd{input}, \cmd{include}, \cmd{includegraphics} (or rather 
the commands specified in the \TeX\ mode variable \acmd{boxMacroNames} 
mentioned in section~\ref{sec:variables}), \cmd{bibliography}, 
\cmd{usepackage}, and \cmd{documentclass}.  Note that the required 
arguments of these commands are files, and so command-double-clicking 
such an argument opens the corresponding file.  Unless the filename 
includes a Macintosh path (which is not recommended, since it's not 
portable), the current folder is searched first.  If the file is not 
found in the current folder, the algorithm next checks the hierarchy 
of folders under the user-specified ``TeX Inputs Folder'', which is 
optionally set by choosing \acmd{App Paths} on the \amenu{Config} menu.  
If the file is still not found, all folders whose name contains the 
string ``inputs'' in the \TeX\ application folder are checked next.  
For example, all folders in the \TeX\ folder with names such as 
\filespec{TeX-inputs}, \filespec{TeX-inputs2}, and 
\filespec{My-TeX-inputs} will be searched.

\textbf{Tip:} A command-double-click operation may be simulated with a 
keystroke.  With the cursor inside the required argument of a 
cmd-dbl-clickable \LaTeX\ command, press \key{F6} to activate the 
algorithm.

\Alpha\ has another modified double-click that will be of interest 
to \LaTeX\ users.  You may already know that double-clicking a 
delimiter (parenthesis, bracket, or brace) selects the text between 
it and its matching delimiter.  Moreover, if you hold down the 
\key{Ctl} key while double-clicking a delimiter, the text \emph{and} 
the delimiters will be selected.  These commands are very handy for 
cutting and pasting blocks of delimited text, especially in a 
\LaTeX\ document where braces, for example, run rampant.

\chapter{Technical tips}

\textbf{Tip:} To see what \TeX\ applications are currently supported 
for typesetting, viewing, or printing, type
\begin{verbatim}
  array names texAppSignatures
\end{verbatim}
or
\begin{verbatim}
  array names viewDVIAppSignatures
\end{verbatim}
or
\begin{verbatim}
  array names printDVIAppSignatures
\end{verbatim}
respectively, in the Tcl shell.  (To invoke the shell, choose the 
\acmd{Shell} command from \Alpha's \amenu{File} menu or press 
\key{Cmd Y}.)

\textbf{Tip:} To see what applications are currently supported for 
creating, viewing, or printing \filetype{ps} files, type
\begin{verbatim}
  array names dvipsAppSignatures
\end{verbatim}
or
\begin{verbatim}
  array names viewPSAppSignatures
\end{verbatim}
or
\begin{verbatim}
  array names printPSAppSignatures
\end{verbatim}
respectively, in the Tcl shell.

\textbf{Tip:} To see what \BibTeX\ or \MakeIndex\ applications are 
currently supported, type
\begin{verbatim}
  array names bibtexAppSignatures
\end{verbatim}
or
\begin{verbatim}
  array names makeindexAppSignatures
\end{verbatim}
respectively, in the Tcl shell.

\end{document}
%%%%%%%%%%%%%%%%%%%%%%%%%%%%%%%%%%%%%%%%%%%%%%%%%%%%%%%%%%%%%%%%%%%%%%%%%%%
